\documentclass[11pt]{article}

\usepackage[a4paper,margin=1in]{geometry}
\usepackage{amsmath,amssymb,amsthm,mathtools}
\usepackage{booktabs}
\usepackage{enumitem}
\usepackage[hidelinks]{hyperref}
\usepackage{xurl}
\usepackage{seqsplit}
\usepackage{microtype}
\usepackage{xcolor}

\newtheorem{theorem}{Theorem}[section]
\newtheorem{lemma}[theorem]{Lemma}
\newtheorem{proposition}[theorem]{Proposition}
\newtheorem{corollary}[theorem]{Corollary}
\newtheorem{remark}[theorem]{Remark}
\newcommand{\R}{\mathbb{R}}
\newcommand{\Ftwo}{\mathbb{F}_2}
\newcommand{\norm}[1]{\left\lVert#1\right\rVert}
\newcommand{\LC}{\mathrm{LC}}
\newcommand{\hash}[1]{\texttt{\seqsplit{#1}}}

\title{Proof-Carrying Finite Planar Three-Body Continuation\\
with Repeated Levi--Civita Chart Switching}
\author{Jeffrey Spaulding\\
\small Independent Researcher\\
\small \href{mailto:jspaulding42@hotmail.com}{jspaulding42@hotmail.com}}
\date{Review draft, \today}

\begin{document}
\maketitle

\begin{abstract}
We present an implemented, proof-carrying continuation checker for one
supplied planar Newtonian three-body point initial-value problem.  Every
accepted raw real numeric field is a finite binary64 value interpreted, in
theorem-facing exact calculations, as its exact dyadic rational.  A raw
version-1 certificate has the finite grammar
\[
 N\bigl(N\mathbin{\to}N\mid
 N\mathbin{\to}\LC_{ij}\mathbin{\to}N\bigr)^*,
 \qquad (i,j)\in\{(0,1),(0,2),(1,2)\},
\]
so all three canonical binary pairs may occur and a pair may be revisited
after returning to ordinary coordinates.  The version-2 replay checker
revalidates every ordinary and Levi--Civita (LC) tube, complete handoff box,
local deck-gauge relation, and conditional physical-clock update.  LC mass
coefficients on this theorem surface are derived from exact mass fractions
and enclosed outward in binary64 arithmetic.  A successful finite induction
produces an exact-clock preimage of a requested physical time and an exact
\texttt{RationalInterval} Horner enclosure of the state there.  The two
mathematical outcomes are \texttt{CERTIFIED\_TO\_T} and a structured
\texttt{UNRESOLVED} frontier; malformed wire data may instead fail before
replay.  The result is replayed by the same implementation and is not an
independent formal verification.  We also retain a separate, stronger local
theorem for a chart with an exact binary-collision event anchor.  An ordinary
LC passage does not by itself prove that a collision occurs.  The contribution
is a sound finite continuation kernel for a specified exact planar IVP, not a
complete producer, a closed form, or an all-time solution of the three-body
problem.
\end{abstract}

\section{Introduction}

The Newtonian three-body flow is classical away from collision, but direct
Cartesian coordinates become singular when two bodies meet.  In the planar
problem the Levi--Civita transformation replaces a selected binary collision
by an analytic lifted differential equation
\cite{LeviCivita1920,SiegelMoser1995,Montgomery2024}.  Regularization alone,
however, does not show that a computed trajectory encloses an exact solution,
that successive charts describe the same initial-value problem, or that their
local time coordinates form one physical clock.  These are proof obligations,
not matters of plotting accuracy.

This paper describes an executable acceptance theorem for a \emph{supplied
finite chain}.  Candidate generation is deliberately outside the theorem.
The checker starts with one exact serialized point IVP, freshly validates the
initial ordinary Newtonian chart, and folds a strict sequence of ordinary
bridges and complete ordinary--LC--ordinary passages.  At every accepted
vertex it carries one actual solution branch and an interval containing that
branch's actual clock origin.  It never promotes all points of a rectangular
tube or all clocks in an interval to realizable trajectories.  Complete box
containment is used to retain the one actual state despite discarded
correlations.

The main contributions are:
\begin{enumerate}[label=(\roman*)]
\item a strict raw certificate grammar whose records do not contain trusted
      intermediate checker results, propagated clocks, gauge choices, or
      final enclosures;
\item a finite conditional-induction theorem for ordinary bridges and for LC
      passages of any canonical binary pair, including repeated visits;
\item exact forward clock-ledger rules and exact fixed-physical-time
      evaluation after the final ordinary chart;
\item an outward mass-coefficient kernel which supports generic positive
      binary64 mass triples on the new LC-tube/carried-entry/carried-exit
      surface without requiring every derived ratio to be exactly binary64;
      and
\item fail-closed replay with either a certified terminal enclosure or a typed
      unresolved frontier retaining only the proved prefix.
\end{enumerate}

The design follows the established validated-integration viewpoint
\cite{Lohner1987,NedialkovJacksonCorliss1999,MooreKearfottCloud2009}, but it
focuses on composition: exact IVP identity, complete transition containment,
chart gauge, and the physical-time cocycle are part of the checked artifact.
Capi\'nski, Kepley, and Mireles James proved substantially stronger
orbit-level collision and near-collision conclusions in the circular
restricted three-body problem using computer-assisted methods
\cite{CapinskiKepleyMirelesJames2023}.  That is a different dynamical model
and proof goal.  The present result concerns one supplied finite chain in the
full planar three-body problem and remains conditional on the trusted kernel
listed in Section~\ref{sec:trusted}.  No novelty or priority relative to that
work is asserted.

The term ``computational solution'' is used here in a deliberately limited
sense: if a finite certificate is accepted, the checker encloses the exact
state at one requested finite time.  A general certified computational
solution would additionally need a producer and a theorem establishing
coverage, progress, and suitable termination for a broad input class.  Those
properties are not consequences of the acceptance theorem proved here.

\section{Ordinary and regularized dynamics}
\label{sec:dynamics}

\subsection{Planar Newtonian charts}

Let the positive masses be $m_0,m_1,m_2$.  For positions
$q_a\in\R^2$ and velocities $v_a\in\R^2$, $a\in\{0,1,2\}$, the ordinary
state $x=(q_0,q_1,q_2,v_0,v_1,v_2)\in\R^{12}$ obeys
\begin{equation}
 \dot q_a=v_a,
 \qquad
 \dot v_a=\sum_{b\ne a}m_b
       \frac{q_b-q_a}{|q_b-q_a|^3}.
 \label{eq:newton}
\end{equation}
An ordinary chart uses a local parameter $s$.  Along the carried physical
solution its physical time has the affine form
\begin{equation}
 t=s+b,
 \label{eq:ordinary-clock}
\end{equation}
where $b$ is the chart's clock origin.  The field is analytic and locally
unique on any box for which every pair distance has a positive lower bound.

\subsection{A selected Levi--Civita chart}

Choose a canonical ordered pair $(i,j)$, with $i<j$, and let $k$ be the
remaining body.  Put
\[
 M=m_i+m_j,\qquad \alpha=\frac{m_j}{M},\qquad
 \beta=\frac{m_i}{M}.
\]
For $z=(z_1,z_2)$ define
\[
 Q(z)=(z_1^2-z_2^2,2z_1z_2),\qquad
 \rho=|z|^2,\qquad
 L(z)=2\begin{pmatrix}z_1&-z_2\\z_2&z_1\end{pmatrix}.
\]
The lifted state is
\[
 X=(z,w,h,R,U,y,V,t)\in\R^{14}.
\]
Here $q=Q(z)$ is the selected relative position, $R$ is its pair center of
mass, and $y$ is the displacement of the third body from $R$.  Define
\[
 d_i=y+\alpha Q(z),\qquad d_j=y-\beta Q(z),
 \qquad f_i=\frac{d_i}{|d_i|^3},\quad
 f_j=\frac{d_j}{|d_j|^3},
\]
\[
 B=\frac{m_k}{M}(m_i f_i+m_j f_j),\qquad
 A_k=-m_i f_i-m_j f_j,
\]
\[
 Y=A_k-B,\qquad P=m_k(f_j-f_i).
\]
The regularized field used by the checker is
\begin{align}
 z'&=w,&
 w'&=\tfrac12hz+\tfrac14\rho L(z)^TP,\nonumber\\
 h'&=(L(z)w)\mathbin{\cdot}P,&
 R'&=\rho U,& U'&=\rho B,\label{eq:lcfield}\\
 y'&=\rho V,& V'&=\rho Y,& t'&=\rho.\nonumber
\end{align}
Primes denote the LC parameter.  There is no division by $z$ or by $\rho$;
the field is analytic at a selected binary collision while $d_i$ and $d_j$
remain separated from zero.

On a punctured LC interval, the physical relative velocity is
\begin{equation}
 v=\frac{L(z)w}{\rho}.
 \label{eq:lc-velocity}
\end{equation}
The Cartesian reconstruction is
\[
 q_i=R-\alpha q,\qquad q_j=R+\beta q,\qquad q_k=R+y,
\]
with the analogous velocity formulas using $U$, $v$, and $V$.  The deck
action
\begin{equation}
 \Gamma(z,w,h,R,U,y,V,t)=(-z,-w,h,R,U,y,V,t)
 \label{eq:deck}
\end{equation}
does not change the physical projection.

The pair-energy constraint is
\begin{equation}
 C=2|w|^2-M-\rho h.
 \label{eq:constraint}
\end{equation}

\begin{lemma}[Constraint and punctured projection]
\label{lem:projection}
For the exact field~\eqref{eq:lcfield}, $C'=0$.  If $C=0$ and $\rho>0$,
the projection~\eqref{eq:lc-velocity} with $dt/ds=\rho$ satisfies the
Newtonian equations~\eqref{eq:newton}.  Moreover
$F(\Gamma X)=D\Gamma F(X)$.
\end{lemma}

\begin{proof}
Since $\rho'=2z\cdot w$, substitution in~\eqref{eq:lcfield} gives
\[
 C'=4w\cdot w'-\rho'h-\rho h'=0.
\]
The identities
\[
 L^TL=LL^T=4\rho I,\qquad Lz=2Q(z),
\]
and
\[
 \rho L(w)w-\rho' L(z)w=-2|w|^2Q(z)
\]
reduce the projected relative acceleration to
\[
 q_{tt}=\frac{(-2|w|^2+\rho h)q}{\rho^3}+P
       =-\frac{Mq}{|q|^3}+P.
\]
The equations for $R,U,y,V$ reconstruct the remaining Newtonian components.
Deck equivariance follows by direct substitution.  The implementation also
contains an exact symbolic audit of these projection and gauge identities.
\end{proof}

\section{A posteriori tubes and arithmetic semantics}
\label{sec:tubes}

Let $F$ denote either the ordinary field or the selected LC field, and let
$\bar X$ be a polynomial on an interval whose maximum distance from its
anchor is $\ell$.  Suppose outward interval evaluation proves, throughout a
declared radius-$r$ tube,
\[
 \norm{\bar X'-F(\bar X)}_\infty\leq\delta,
 \qquad \norm{DF}_\infty\leq K,
\]
and the exact anchored state is within $\varepsilon$ of the polynomial
anchor.  Set
\begin{equation}
 E=e^{K\ell}\varepsilon+
   \delta\frac{e^{K\ell}-1}{K},
 \label{eq:gronwall}
\end{equation}
with $E=\varepsilon+\delta\ell$ when $K=0$.

\begin{lemma}[A posteriori tube]
\label{lem:tube}
If every denominator of $F$ is separated from zero on the declared tube and
$E<r$, then the exact anchored IVP has a unique solution on the complete
chart interval and $\norm{X-\bar X}_\infty\le E$ there.
\end{lemma}

\begin{proof}
Subtract the exact and polynomial integral equations.  Gr\"onwall's
inequality yields~\eqref{eq:gronwall} up to any first exit.  At a hypothetical
first exit the strict estimate gives $E<r$, a contradiction.  Local
uniqueness joins the anchored solution across the interval.
\end{proof}

The ordinary checker uses an anisotropic version of this lemma.  With
position and velocity radii $R_q,R_v$, it evaluates the norm
\[
 \|(\Delta q,\Delta v)\|_W=
 \max\{\|\Delta q\|_\infty/R_q,
       \|\Delta v\|_\infty/R_v\}.
\]
If $A$ bounds the acceleration Jacobian, a valid Lipschitz bound is
\[
 L_W=\max\{R_v/R_q,\,A R_q/R_v\}.
\]
The corresponding weighted defects and anchor errors are inserted in
Lemma~\ref{lem:tube}.  This avoids forcing the position and velocity errors
to have the same scale near a binary encounter.

\subsection{Exact binary64-dyadic inputs}

The raw version-1 loader accepts finite built-in binary64 values in real
numeric fields; schema versions and counts are strict built-in integers.  In
every theorem-facing rational calculation, a binary64 value is interpreted
as the exact dyadic rational represented by its bit pattern.  Thus the token
\texttt{0.1} denotes the exact round-tripped binary64 number, not the
mathematical rational $1/10$.  Exact equality means equality of the
corresponding rational values.  This convention prevents a certificate from
silently changing the serialized IVP to a nearby decimal problem.

Exact integer and \texttt{Fraction} arithmetic is used for polynomial
endpoints and handoff translations.  Other algebraic and transcendental
bounds required by the tube checkers use the directed outward operations
described in Section~\ref{sec:trusted}.

\subsection{Version-2 outward mass coefficients}

Generic binary64 masses need not have binary64-representable derived ratios.
The new theorem-facing LC tube, carried entry, and carried exit therefore do
not compare a rounded ratio with an exact ratio.  They first convert
$m_i,m_j,m_k$ to exact fractions and derive exactly
\begin{equation}
 M,\quad \frac{m_j}{M},\quad \frac{m_i}{M},\quad
 \frac{m_k}{M},\quad
 \frac{m_i m_k}{M},\quad \frac{m_j m_k}{M},\quad
 m_i\left(1+\frac{m_k}{M}\right),\quad
 m_j\left(1+\frac{m_k}{M}\right).
 \label{eq:mass-coefficients}
\end{equation}
Each exact coefficient is then converted to the tight enclosing binary64
interval before it is consumed by the interval-dual LC field, its Jacobian,
or the complete projection formulas.

\begin{proposition}[Mass-system fidelity]
\label{prop:masses}
Subject to sound directed conversion and interval arithmetic, every accepted
version-2 LC tube, carried entry, and carried exit encloses the exact LC field
and projection for the exact binary64-dyadic mass triple serialized at the
root.
\end{proposition}

\begin{proof}
All coefficients in~\eqref{eq:mass-coefficients} are rational functions of
positive exact rational masses, so their exact values are defined.  Inclusion
isotonic evaluation with an outward interval containing each value encloses
the same exact-mass field and projection.  No nearby rounded mass tuple is
substituted.
\end{proof}

An obligation may remain unresolved if an outward coefficient is not finite
or if a later interval operation overflows.  That is checker incompleteness,
not a statement about the physical solution.  This migration is intentionally
limited to the v2 LC-tube/carried-entry/carried-exit/finite-chain surface;
legacy LC checking paths retain their documented arithmetic gates.

\section{The raw finite-chain certificate}
\label{sec:grammar}

The wire schema is \texttt{RawPlanarChainCertificate} version~1.  The current
top-level replay checker is version~2.  The root contains positive masses,
three initially distinct planar positions, three velocities, an initial
physical time, and an ordinary-chart parameter, together with one ordinary
chart and its tube.  Acceptance requires zero binding tolerances and exact
agreement among the binding parameter, the initial chart's left endpoint,
and the tube anchor.  The initial chart and tube are freshly replayed.  The
induction therefore starts from one exact point IVP rather than an
initial-condition box.

The strict tagged-union grammar is
\begin{equation}
 N\bigl(N\to N\mid N\to\LC_{ij}\to N\bigr)^*,
 \qquad (i,j)\in\{(0,1),(0,2),(1,2)\}.
 \label{eq:grammar}
\end{equation}
An \texttt{ordinary\_bridge\_v1} record contains a raw ordinary transition,
target ordinary chart, and target tube.  A
\texttt{planar\_lc\_passage\_v1} record contains a raw ordinary-to-LC entry,
LC chart and tube, LC-to-ordinary exit, target ordinary chart, and target
tube.  Every LC segment returns to $N$ before the next segment begins.
Consequently the three pair charts can occur in any finite order, including
revisits, without a direct transition between different LC coordinate
systems.

Parsing is fail-closed.  Exact classes, field sets, tags, canonical pair
labels, round trips, and a globally unique identifier namespace are required.
The raw wire does not contain a later IVP binding, a propagated clock origin,
a selected gauge assignment, a projected endpoint, a cached checker result,
or a final enclosure.  Those objects are derived during replay.  Canonical
serialization and its aggregate SHA-256 bind transport and mutation of the
raw record; the digest is not a substitute for any mathematical obligation.

\section{The conditional induction kernel}
\label{sec:induction}

After the root and after every fully accepted segment, the current vertex is
an ordinary chart.  The proof carries the following invariant.

\begin{proposition}[Carried-branch invariant]
\label{prop:invariant}
At a certified ordinary vertex there is one actual branch of the unique
solution of the exact root IVP inside the freshly checked current tube.  Its
ordinary clock is $t=s+b$ for one actual scalar $b$ in a checker-derived exact
rational interval $B=[B_-,B_+]$.  All accepted handoffs identify this branch
with the same physical solution, modulo a pair-local LC deck transformation.
\end{proposition}

This statement is existential in both the branch and the clock.  It does not
say that every state in the tube is on the root trajectory, that every
$b\in B$ is realizable, or that arbitrary combinations of state and clock
interval coordinates are correlated.  Complete interval containment is used
only to prove that the one actual carried state remains included.

At the root, if the serialized initial physical time is $t_0$ and the initial
ordinary parameter is $s_0$, the invariant holds with the exact point
\begin{equation}
 B_0=[t_0-s_0,t_0-s_0].
 \label{eq:root-clock}
\end{equation}

\subsection{Ordinary bridge}

Let $e$ be the source ordinary chart's right parameter and $a$ the target
ordinary chart's left parameter.  The bridge helper freshly checks both
tubes, evaluates every source polynomial over the complete endpoint by exact
rational arithmetic, inflates by the checked source error, and requires the
resulting twelve-dimensional box to lie inside the target initial ball.
Center or sampled agreement is insufficient.

For the actual carried branch, ordinary autonomy and uniqueness then identify
the target solution with the source solution.  Since the physical endpoint
time is $e+b$, the target clock origin is $b'=e+b-a$.  The interval update is
derived exactly:
\begin{equation}
 B'=B+e-a.
 \label{eq:n-to-n-clock}
\end{equation}
The cocycle record stores both exact intervals and the exact parameter
translation; no producer-supplied clock is accepted.

\subsection{Complete ordinary-to-LC entry}

At the ordinary right endpoint the entry helper reconstructs the complete
source box and requires the selected pair to be separated throughout it.  A
deterministic square-root grammar produces a complete finite LC lift cover.
Depending on its relation to the square-root branch cut, that cover contains
one or two patches.  In the two-patch case their overlap determines an exact
parity-one relation under~\eqref{eq:deck}.  The checker constructs the exact
$\Ftwo$ graph, enumerates its two global complements, and accepts only if one
coherent complement places every complete emitted patch box inside one LC
target anchor.

If $D_{\rm in}$ denotes the interval of entry physical times, it is derived
from the source ledger as
\begin{equation}
 D_{\rm in}=B+e.
 \label{eq:entry-clock}
\end{equation}
For the one actual carried source state, patch coverage supplies at least one
exact constrained LC lift.  The accepted coherent gauge sends that lift into
the checked LC anchor, and its physical-time coordinate lies in
$D_{\rm in}$.  The claim is not that every point of every rectangular patch
box satisfies the nonlinear pair-energy constraint.

\subsection{LC tube and strict physical time}

The v2 LC checker freshly validates the entire fourteen-dimensional tube,
including residual, derivative, Gr\"onwall, third-body separation, and mass
arithmetic obligations.  The carried entry satisfies~\eqref{eq:constraint},
Lemma~\ref{lem:projection} preserves it, and the punctured projection is the
same exact-mass Newtonian branch.

The entry is separated, so $\rho>0$ there and the analytic function $z$ is not
identically zero.  Its zeros are therefore isolated.  Since
\[
 \frac{dt}{ds}=|z|^2\ge0,
\]
the physical-time increment over every nondegenerate LC parameter interval is
strictly positive.  This proves correct time ordering even if an isolated
zero occurs.  It does not prove that $z$ has a zero.

\subsection{Complete LC-to-ordinary exit}

At the exact LC right parameter the exit helper evaluates all fourteen
components and inflates them by the freshly checked LC error.  It requires a
strict positive lower bound for $\rho$ on the complete exit slice.  The full
slice is then projected with the outward mass coefficients to twelve
Cartesian position and velocity intervals, all of which must be contained in
the target ordinary initial ball.

The LC physical-time component gives an outward interval
\[
 D_{\rm out}=[d_-,d_+]\ni t_{\rm exit}.
\]
For target ordinary anchor $a$, the checker derives
\begin{equation}
 B'=D_{\rm out}-a.
 \label{eq:lc-to-n-clock}
\end{equation}
Constraint invariance, strict LC time, deck-equivariant projection, complete
endpoint containment, and ordinary uniqueness reestablish
Proposition~\ref{prop:invariant} at the target vertex.

\subsection{Local gauges and repeated pairs}

Each LC passage has its own patch namespace and one coherent gauge assignment.
Gauges belonging to different selected pairs are never compared.  Two visits
to the same pair separated by an ordinary chart are also independent lift
problems; the checker does not impose an unsupported equality between their
gauge bits.  Ordinary Cartesian coordinates are the common physical
representation through which pair changes and revisits compose.

\section{Fixed-time soundness theorem}
\label{sec:main-theorem}

After all supplied segments are accepted, let the current ordinary chart have
domain $[a,r]$, and let the carried clock origin satisfy
$b\in B=[B_-,B_+]$.  For a finite requested physical target $T$, the checker
first requires
\begin{equation}
 T\ge a+B_+.
 \label{eq:target-forward}
\end{equation}
It then forms the exact parameter preimage interval
\begin{equation}
 J=T-B=[T-B_+,T-B_-]
 \label{eq:target-preimage}
\end{equation}
and requires $J\subseteq[a,r]$.  Each final position and velocity polynomial
is evaluated on all of $J$ by exact \texttt{RationalInterval} Horner
arithmetic.  The current ordinary tube is replayed and its validated error is
added outward.  Exact rational endpoint differences determine the largest
component width, which must not exceed the requested bound.

\begin{theorem}[Raw finite planar continuation-chain soundness]
\label{thm:chain}
Let $C$ be a canonically constructed raw version-1 certificate whose root is
one exact positive-mass planar point IVP with initially distinct positions,
and let $T$ be its finite target.  Subject to the trusted kernel of
Section~\ref{sec:trusted}, if version-2 replay returns
\[
 \operatorname{check}(C).\mathrm{status}
 =\texttt{CERTIFIED\_TO\_T},
\]
then:
\begin{enumerate}[label=(\roman*)]
\item the exact root IVP has one unique chain-compatible generalized planar
      continuation through the supplied finite chain to physical time $T$;
\item its ordinary pieces solve~\eqref{eq:newton}, and every punctured
      constrained LC projection is that same Newtonian solution;
\item every handoff carries the same actual branch at one globally coherent
      physical clock, with strictly forward time through each LC passage;
\item all selected pair changes and repeated visits compose through ordinary
      coordinates with only pair-local deck ambiguity;
\item the returned ordinary box contains the state at exactly $T$; and
\item every returned component width satisfies the requested bound.
\end{enumerate}
Uniqueness here is relative to the supplied chain and its selected regularized
lift; it is not a classification of every possible generalized continuation
at a collision.
\end{theorem}

\begin{proof}
The exact root binding and fresh ordinary tube establish
Proposition~\ref{prop:invariant} with~\eqref{eq:root-clock}.  Induct on the
finite segment tuple.  A complete ordinary bridge preserves the actual branch
and derives~\eqref{eq:n-to-n-clock}.  A complete LC passage covers the actual
ordinary endpoint by a constrained lift, selects one coherent local gauge,
validates its unique lifted evolution, projects the full separated exit into
the next ordinary anchor, and derives~\eqref{eq:lc-to-n-clock}.  Thus every
accepted segment reestablishes the invariant for the same root solution.
LC strict-clock monotonicity forbids time reversal.  At the terminal vertex,
the actual $b\in B$ implies $s_T=T-b\in J$.  Exact Horner evaluation on all of
$J$, followed by the fresh tube inflation, contains the state at $T$.
The exact width gate proves the last conclusion.
\end{proof}

\begin{remark}[No event conclusion from an ordinary LC passage]
Theorem~\ref{thm:chain} permits regularized coordinate passages.  It does not
assert that $z=0$, or equivalently $\rho=0$, occurs in any such passage.
Collision occurrence, isolation, and time require a separate checked event
witness.
\end{remark}

\section{Unresolved frontiers and replay semantics}
\label{sec:unresolved}

Malformed canonical JSON, unknown fields or tags, wrong exact classes, and
other parser violations may raise before any replay result exists.  For a
successfully constructed exact-class certificate, a failed tube,
containment, arithmetic, target-domain, or width obligation returns
\texttt{UNRESOLVED}.  It is not a claim of nonexistence, blow-up, or collision.

Replay counts the longest consecutive segment prefix accepted before the
first failed segment and retains at most one typed region:
\begin{itemize}
\item the current ordinary right frontier after a certified root or prefix;
\item a fourteen-dimensional LC right frontier when entry and LC-tube replay
      succeed but the failing passage's exit does not; or
\item the ordinary fixed-$T$ enclosure when mathematical enclosure succeeds
      but the requested width gate fails.
\end{itemize}
No certified root means no certified prefix and no retained region.  An LC
frontier is never relabeled as a finite Cartesian velocity state.  If a
frontier time interval contains the unknown actual endpoint, only its lower
endpoint is reported as guaranteed reached.  An unresolved result proves only
the accepted prefix and its typed containment statement; it says nothing
about continuation or certifiability beyond that frontier.

The result retains the raw certificate, its canonical aggregate digest,
derived ledgers, checker and kernel identifiers, and the obligation ledger.
Its certification property reruns the same top-level implementation and
requires exact result equality.  Mutated evidence, hostile subclasses,
copied all-true ledgers, changed identifiers, and altered derived snapshots do
not inherit certification.  This is exact-self replay, not an independently
implemented verifier or a proof-assistant proof.

For review, a successful top-level result pins the identifiers
\path{raw_planar_continuation_chain_replay_checker_v2},
\path{exact_root_and_conditional_planar_segment_induction_kernel_v1},
\path{forward_interval_clock_origin_ledger_v1}, and
\path{exact_rational_horner_fixed_time_kernel_v1}.  Its segment ledger admits
only \path{carried_ordinary_bridge_checker_v1} and
\path{carried_planar_lc_exit_checker_v2}.  The latter self-binds
\path{carried_planar_lc_entry_checker_v2},
\path{independent_planar_lc_aposteriori_tube_checker_v2}, and the mass kernel
\path{planar_lc_mass_coefficients_exact_binary64_fraction_outward_v1}.
These strings are implementation identity pins, not evidence that the
so-named components are independent implementations of the top-level replay.

\section{A separate event-anchored collision theorem}
\label{sec:collision}

The repository retains a stronger local checker for a supplied LC chart whose
anchor is asserted to be the collision event.  It is logically separate from
the generic finite-chain passage.

\begin{theorem}[Exact event-anchored binary-collision passage]
\label{thm:collision}
Fix positive masses and one selected pair.  Suppose the event checker accepts
an LC tube on $[s_-,s_+]$, $s_-<0<s_+$, with zero anchor error at $s=0$, and
checks exactly
\begin{equation}
 z(0)=0,\qquad 2|w(0)|^2=M.
 \label{eq:collision-anchor}
\end{equation}
Suppose it also validates separated-third-body denominators on the complete
LC tube, positive $\rho$ on both complete punctured endpoint slices, complete
LC-to-Newton projection there, and containment in validated ordinary tubes.
Then the exact lifted IVP has an isolated selected-pair collision at $s=0$;
its two punctured projections are branches of the same Newtonian trajectory,
and physical time is strictly ordered from the left endpoint through the
collision to the right endpoint.
\end{theorem}

\begin{proof}
Lemma~\ref{lem:tube} gives one analytic lifted solution.  The exact anchor
sets $C=0$ and, because $M>0$, gives $w(0)\ne0$.  Hence $z$ has a simple
isolated zero.  In particular,
\[
 t(s)-t(0)=\int_0^s|z(u)|^2\,du
 =\frac{|w(0)|^2}{3}s^3+O(s^4).
\]
Lemma~\ref{lem:projection} supplies the Newtonian equations off the event.
Projection of the complete endpoint slices and ordinary uniqueness bind both
punctured branches to the same lifted IVP.  Finally, $t'=|z|^2$ and
nonidentity of analytic $z$ imply strict physical-time ordering.
\end{proof}

The zero-error constraint anchor is essential: a positive box centered on a
constrained point generally includes off-constraint states, whose projected
singular coefficient is not the same Newtonian coefficient.  This theorem is
therefore an event theorem only when the exact identities
\eqref{eq:collision-anchor} are checked.  The retained event checker has its
own documented mass-consistency restrictions; Proposition~\ref{prop:masses}
does not claim that every legacy collision or overlap path has been migrated
to the v2 outward mass kernel.

\section{Review artifact and reproducibility}
\label{sec:artifact}

The primary review evidence is raw JSON, not the candidate builder.  The
script
\begin{center}
\path{scripts/certify_repeated_planar_chain.py}
\end{center}
constructs candidate fixtures, exports canonical raw records, performs fresh
replay, and verifies a transport manifest.  Its candidate-construction path
is untrusted: only the raw record and the obligations recomputed from that
record enter Theorem~\ref{thm:chain}.

The version-0.3 review bundle is located at
\begin{center}
\path{artifacts/v0.3.0-review/planar-chain/}
\end{center}
and is organized around the following files:
\begin{itemize}
\item \path{success.raw.json} and \path{success.replay.json};
\item \path{failed-revisit.raw.json} and
      \path{failed-revisit.replay.json}; and
\item \path{manifest.json}.
\end{itemize}
The success fixture has the five-segment chart word
\[
 N\to N\to\LC_{01}\to N\to\LC_{02}\to N
 \to\LC_{12}\to N\to\LC_{01}\to N.
\]
It exercises all canonical pairs and a same-pair revisit.  The companion
negative fixture is designed so that the final revisit exit fails after a
four-segment certified prefix, retaining a typed fourteen-dimensional LC
frontier.  These expected fixture outcomes are regression evidence for the
finite grammar and failure semantics, not proof that any LC passage contains
a physical collision.

For the candidate bytes selected for the review bundle, fresh
same-implementation replay gave the following checked values.  The success
raw SHA-256 is
\hash{ede15b0f35cf741f542a6cd260470a93ee5ff85dc5ae2371db1b88819b821f11};
its status is \texttt{CERTIFIED\_TO\_T} after five segments.  The requested
target is approximately $0.008008225523335941$, the exact final parameter
preimage converts to approximately
$[0.005,0.009000009750858424]$, and the reported maximum component width is
approximately $0.021640584283380553<0.1$.  The negative raw SHA-256 is
\hash{c6830919726c22fbca6e45d5781b2465e54876bfa89c25bc26e97284ff918727};
its status is \texttt{UNRESOLVED}, with prefix count four, failed segment
index four, first failing obligation
\path{carried_lc_exit_target_initial_ball_contains_complete_projection}, and
a retained fourteen-dimensional LC frontier.

The canonical manifest remains the source of the machine-readable values.
Its hashes bind transport bytes, not mathematical truth.  Environment
versions are informational and are not silently promoted to theorem
assumptions.  Fresh-process export comparison and bundle verification are
used to detect nondeterministic serialization.  The replay transcripts are
generated by the same Python implementation as the checker and must not be
described as an independent verification.

The local projection algebra audit remains available through
\begin{center}
\path{scripts/verify_lc_projection_identities.py}.
\end{center}
The repository documentation records focused and full regression commands.
Release review should use the tracked raw bytes, rerun the checker in a clean
environment, compare the new transcript with the tracked one, and report any
difference rather than editing a result record by hand.

\section{Trusted kernel and precise nonclaims}
\label{sec:trusted}

Theorem~\ref{thm:chain} is conditional on the following trusted kernel:
\begin{itemize}
\item Python integer and \texttt{Fraction} arithmetic, strict class checks,
      canonical JSON handling, and SHA-256 transport hashing;
\item exact \texttt{RationalInterval} polynomial arithmetic and
      binary64-to-exact-fraction conversion;
\item directed \texttt{nextafter} inflation and the scalar outward bounds
      used for square roots, reciprocal powers, exponentials, Gr\"onwall
      estimates, square-root lifts, and projections;
\item interval forward automatic differentiation for the ordinary and LC
      fields, including complete Jacobian row-sum bounds;
\item the explicit Newtonian and LC vector fields, mass-coefficient formulas,
      constraint, lift, projection, and deck action;
\item the a posteriori ODE lemma, analytic local existence and uniqueness,
      constraint invariance, punctured projection equivalence, deck
      equivariance, and the nontrivial-analytic-$z$ strict-clock lemma;
\item complete interval containment, the canonical square-root cover, the
      exact $\Ftwo$ graph kernel, conditional segment induction, the forward
      clock ledger, and fixed-time evaluation; and
\item correctness of the Python runtime and of the implementation of these
      kernels.
\end{itemize}
IEEE~1788 provides historical standardization context for interval arithmetic
\cite{IEEE1788_2015}; a general reference is
Moore, Kearfott, and Cloud~\cite{MooreKearfottCloud2009}.  In the present
binary64 backend, sensitive quantities are inflated outward and some scalar
bounds use directed high-precision decimal arithmetic.  In particular,
\texttt{Decimal.exp} under its directed context is part of the runtime
hypothesis; this manuscript does not prove it to be a correctly rounded
transcendental implementation.  Adversarial tests support the arithmetic
kernel but do not replace its audit.

The current result does \emph{not} establish any of the following:
\begin{itemize}
\item a complete or terminating certificate producer for arbitrary initial
      data;
\item a theorem for boxes or families of IVPs;
\item that every point in a state or clock interval is realizable;
\item automatic collision detection, collision occurrence in an ordinary LC
      passage, or collision isolation without a separate event witness;
\item continuation through simultaneous or total collision;
\item spatial motion, Kustaanheimo--Stiefel chart gluing, or a spatial
      finite-chain theorem;
\item escape completeness, non-Zeno progress, all-time coverage, or an
      infinite-atlas theorem;
\item an elementary expression, closed form, or general closed solution of
      the three-body problem;
\item migration of every legacy LC checker to outward mass arithmetic;
\item an independently implemented verifier, an arbitrary-precision audited
      backend, or a proof-assistant formalization;
\item independent external mathematical review or institutional endorsement;
      or
\item correctness without the arithmetic, analytic, runtime, and
      implementation assumptions above.
\end{itemize}

Sampled residuals, nominal physical-time metadata, producer success flags,
reason strings, producer-selected gauge bits, object subclasses, and cached
checker results are not theorem hypotheses.  The scientifically defensible
claim is narrower: accepted raw evidence carries one supplied exact planar
point IVP through one supplied finite ordinary/LC chain to a rigorous
fixed-time enclosure.

\section{Directions toward broader certified computation}

The present theorem supplies the soundness half of a proof-carrying
integrator: candidate construction may be heuristic, but every accepted step
is checked from raw evidence and composes through an explicit invariant.  A
broader computational result now depends less on adding chart types to the
statement and more on proving progress properties for a producer.  Concrete
next steps are:
\begin{enumerate}[label=(\roman*)]
\item build adaptive ordinary and LC certificate producers whose output is
      measured by acceptance rate, certified time advance, and enclosure
      growth rather than nominal numerical success;
\item add rigorous event isolation when a collision conclusion is required,
      instead of inferring an event from LC chart use;
\item construct an independent verifier, preferably with an audited
      arbitrary-precision interval backend, and compare canonical raw
      transcripts across implementations;
\item formalize the finite induction and clock ledger in a proof assistant,
      leaving only explicit arithmetic certificates to an executable kernel;
\item study validated families and parameter boxes, beginning with robust
      noncollision regimes and isolated binary-encounter scattering maps; and
\item formulate and test quantitative progress hypotheses needed for escape
      coverage, collision classification, and exclusion of Zeno switching.
\end{enumerate}
None of these steps is required for the finite acceptance theorem, and none is
claimed by the current artifact.  Together they give a falsifiable path from
one proof-carrying finite trajectory to a genuinely broader certified
computational solution.

\section{Software availability}

The implementation is available under the BSD-3-Clause license at
\begin{center}
\url{https://github.com/jspaulding42/three-body-regularized-atlas}.
\end{center}
The immutable tag \texttt{v0.2.0-review} contains the earlier local
collision-passage snapshot.  This manuscript describes the prospective
\texttt{v0.3.0-review} finite-chain artifact; no v0.3 tag or archival DOI is
claimed until the release is actually published.  The repository contains
the raw schema, replay implementation, candidate builder, tests, review
bundle, and the exact symbolic projection audit.

\section*{Generative-AI use and contributions}

This project was conducted through human-directed use of OpenAI generative-AI
systems identified in the user interface as GPT-5.5, GPT-5.5 Pro, and GPT-5.6.
Those systems contributed substantially to mathematical exploration, software
implementation, certificate and test design, manuscript drafting, and critical
revision.  Jeffrey Spaulding initiated and directed the project, selected its
research direction and scope, reviewed and authorized the release of the
resulting artifact, and is the sole accountable author.  The AI systems are
not authors: they cannot approve the manuscript, hold rights, disclose
conflicts, or accept responsibility for its contents.  The manuscript and
software should therefore be read as a human-released, AI-assisted research
artifact whose mathematical claims remain subject to independent expert
review.  No institutional endorsement or independent external mathematical
review is asserted.

\begin{thebibliography}{99}

\bibitem{CapinskiKepleyMirelesJames2023}
M.~J. Capi\'nski, S.~Kepley, and J.~D. Mireles James,
``Computer assisted proofs for transverse collision and near collision orbits
in the restricted three body problem,''
\emph{Journal of Differential Equations}, vol.~366, pp.~132--191, 2023.
\href{https://doi.org/10.1016/j.jde.2023.03.053}{doi:10.1016/j.jde.2023.03.053}.

\bibitem{IEEE1788_2015}
IEEE Standards Association,
\emph{IEEE Standard for Interval Arithmetic}, IEEE Std 1788-2015, 2015.
\href{https://doi.org/10.1109/IEEESTD.2015.7140721}{doi:10.1109/IEEESTD.2015.7140721}.

\bibitem{LeviCivita1920}
T.~Levi-Civita,
``Sur la r\'egularisation du probl\`eme des trois corps,''
\emph{Acta Mathematica}, vol.~42, pp.~99--144, 1920.
\href{https://doi.org/10.1007/BF02404404}{doi:10.1007/BF02404404}.

\bibitem{Lohner1987}
R.~J. Lohner,
``Enclosing the solutions of ordinary initial and boundary value problems,''
in \emph{Computer Arithmetic: Scientific Computation and Programming
Languages}, E.~Kaucher, U.~Kulisch, and C.~Ullrich, Eds. Stuttgart: Teubner,
1987, pp.~255--286.

\bibitem{Montgomery2024}
R.~Montgomery,
``Regularizing binary collisions,'' in \emph{Four Open Questions for the
$N$-Body Problem}. Cambridge: Cambridge University Press, 2024,
pp.~271--275.
\url{https://doi.org/10.1017/9781009200608.017}.

\bibitem{MooreKearfottCloud2009}
R.~E. Moore, R.~B. Kearfott, and M.~J. Cloud,
\emph{Introduction to Interval Analysis}. Philadelphia: Society for
Industrial and Applied Mathematics, 2009.
\href{https://doi.org/10.1137/1.9780898717716}{doi:10.1137/1.9780898717716}.

\bibitem{NedialkovJacksonCorliss1999}
N.~S. Nedialkov, K.~R. Jackson, and G.~F. Corliss,
``Validated solutions of initial value problems for ordinary differential
equations,'' \emph{Applied Mathematics and Computation}, vol.~105, no.~1,
pp.~21--68, 1999.
\href{https://doi.org/10.1016/S0096-3003(98)10083-8}{doi:10.1016/S0096-3003(98)10083-8}.

\bibitem{SiegelMoser1995}
C.~L. Siegel and J.~K. Moser,
\emph{Lectures on Celestial Mechanics}. Berlin: Springer, 1995 reprint.
\href{https://doi.org/10.1007/978-3-642-87284-6}{doi:10.1007/978-3-642-87284-6}.

\end{thebibliography}

\end{document}
